\documentclass{article}
\usepackage{amsmath}
\usepackage[equations]{propositions}

%% Stray spaces around a name must not reach the label or the reference.
%%
%% l3keys strips spaces around a key's value, which covers `name = Alpha ',
%% but two routes bypass it: a bare name, where the whole argument is wrapped
%% as name={...} verbatim, and a braced value, where the trimming stops at the
%% brace.  Untrimmed, the leading space of \item[ Alpha ] was typeset inside
%% the label box and pushed the label 3.8pt off the margin.

\begin{document}

\section{Every spelling of the same name places identically}

The rule marks the text margin.  All five labels must begin at it.

\noindent\rule{\textwidth}{0.3pt}
\begin{prop}
  \item[Alpha] Bare name. \label{s:1}
  \item[ Alpha ] Bare name written with spaces around it. \label{s:2}
  \item[name=Alpha] Explicit key. \label{s:3}
  \item[name = Alpha ] Explicit key with spaces, which l3keys already
    trimmed. \label{s:4}
  \item[name = { Alpha }] Braced value, where l3keys cannot reach the
    spaces. \label{s:5}
\end{prop}
\noindent\rule{\textwidth}{0.3pt}

The references must match one another too, since the name feeds both the label
and the reference: \ref{s:1}/\ref{s:2}/\ref{s:3}/\ref{s:4}/\ref{s:5}.

\section{A space asked for explicitly survives}

Trimming removes only incidental space.  A name that really does want to begin
with one can say so, and the label below must sit further right than those
above.

\noindent\rule{\textwidth}{0.3pt}
\begin{prop}
  \item[{~Alpha}] Deliberate leading space, written \verb|~|.
\end{prop}
\noindent\rule{\textwidth}{0.3pt}

\section{Equation tags are handed a clean number}

\verb|\tagform@| used to pass the display format
\verb|\ignorespaces #1 \unskip \@@italiccorr|, amsmath's guards for the number,
which became part of \verb|#1|.  A format that prints \verb|#1| never noticed;
one that uses it as data --- an index key, say --- got the guards in the key.
The guards are now applied to the token list instead, so \verb|\tag{ 5 }| is
still tagged \texttt{5} rather than \texttt{\ 5\ }.

\begingroup
\SetPropStyle{equation}{display format = [#1]}
\begin{equation} a = b \end{equation}
\begin{equation} c = d \tag{ 5 } \end{equation}
\begin{equation} e = f \tag*{ 7 } \end{equation}
\endgroup

Expected: \texttt{[1]}, then \texttt{[5]} with the spaces gone, then
\texttt{7} --- \verb|\tag*| bypasses the hook, which is how \verb|\ptag|
reaches the format by the ordinary item path instead.

\end{document}
